\documentclass{beamer}
% =====================================================================
% finance-beamer.tex — shared Beamer style for the LLM-in-Finance course
% =====================================================================
% Reusable slide style. Each deck \input's this immediately after
% \documentclass{beamer}. It provides:
%   * the Madrid theme + book colour palette (matches the printed book)
%   * two book-style framing blocks:
%       \begin{bigpicture} ... \end{bigpicture}   ("the bigger picture")
%       \begin{underhood}   ... \end{underhood}    ("under the hood")
%     These mirror the book's `context` and `deepdive` tcolorboxes so the
%     same intuition-vs-mechanism scaffold the reader meets in the book
%     reappears on the slides.
%   * a Python listings style for code frames
%   * appendix conventions: \appendixstart drops a divider and switches
%     to non-numbered backup frames so "extra / less central" material
%     lives after the main talk without inflating the slide count.
%
% Usage:
%   \documentclass{beamer}
%   % =====================================================================
% finance-beamer.tex — shared Beamer style for the LLM-in-Finance course
% =====================================================================
% Reusable slide style. Each deck \input's this immediately after
% \documentclass{beamer}. It provides:
%   * the Madrid theme + book colour palette (matches the printed book)
%   * two book-style framing blocks:
%       \begin{bigpicture} ... \end{bigpicture}   ("the bigger picture")
%       \begin{underhood}   ... \end{underhood}    ("under the hood")
%     These mirror the book's `context` and `deepdive` tcolorboxes so the
%     same intuition-vs-mechanism scaffold the reader meets in the book
%     reappears on the slides.
%   * a Python listings style for code frames
%   * appendix conventions: \appendixstart drops a divider and switches
%     to non-numbered backup frames so "extra / less central" material
%     lives after the main talk without inflating the slide count.
%
% Usage:
%   \documentclass{beamer}
%   % =====================================================================
% finance-beamer.tex — shared Beamer style for the LLM-in-Finance course
% =====================================================================
% Reusable slide style. Each deck \input's this immediately after
% \documentclass{beamer}. It provides:
%   * the Madrid theme + book colour palette (matches the printed book)
%   * two book-style framing blocks:
%       \begin{bigpicture} ... \end{bigpicture}   ("the bigger picture")
%       \begin{underhood}   ... \end{underhood}    ("under the hood")
%     These mirror the book's `context` and `deepdive` tcolorboxes so the
%     same intuition-vs-mechanism scaffold the reader meets in the book
%     reappears on the slides.
%   * a Python listings style for code frames
%   * appendix conventions: \appendixstart drops a divider and switches
%     to non-numbered backup frames so "extra / less central" material
%     lives after the main talk without inflating the slide count.
%
% Usage:
%   \documentclass{beamer}
%   \input{../_style/finance-beamer.tex}
%   \title{...}\subtitle{...}\author{...}\institute{...}\date{\today}
%   \begin{document} ... \end{document}
% =====================================================================

\usetheme{Madrid}
\usecolortheme{whale}
\setbeamertemplate{navigation symbols}{}
\setbeamertemplate{caption}[numbered]
\setbeamercovered{transparent}

% --- Density controls: keep dense, equation-heavy frames on one slide ---
% Slightly smaller body text + tighter lists/blocks. Frame titles and the
% Madrid chrome keep their normal size, so the deck still reads as Madrid,
% just with more vertical room for content.
\setbeamerfont{normal text}{size=\small}
\setbeamertemplate{itemize/enumerate body begin}{\small}
\setbeamertemplate{itemize/enumerate subbody begin}{\footnotesize}
% Tighter block spacing.
\setbeamertemplate{block begin}{%
  \par\vskip2pt%
  \begin{beamercolorbox}[colsep*=.75ex,leftskip=1ex,rightskip=1ex]{block title}%
    \usebeamerfont*{block title}\insertblocktitle%
  \end{beamercolorbox}%
  {\parskip0pt\vskip-1pt}%
  \begin{beamercolorbox}[colsep*=.75ex,vmode,leftskip=1ex,rightskip=1ex]{block body}%
  \vskip1pt}
\setbeamertemplate{block end}{\end{beamercolorbox}\vskip2pt}

\usepackage{amsmath, amssymb}
\usepackage{booktabs}
\usepackage{xcolor}
\usepackage{listings}
\usepackage[skins, breakable]{tcolorbox}

% --- Book colour palette (kept in sync with book/preamble.tex) ---
\definecolor{linkblue}{RGB}{14, 75, 140}
\definecolor{deepdivebg}{RGB}{235, 244, 255}
\definecolor{narrativebg}{RGB}{255, 252, 240}
\definecolor{narrativerule}{RGB}{180, 130, 40}
\definecolor{steelblue}{RGB}{70, 130, 180}
\definecolor{forestgreen}{RGB}{34, 139, 34}
\definecolor{crimson}{RGB}{220, 20, 60}
\definecolor{darkorange}{RGB}{255, 140, 0}
\definecolor{codebg}{RGB}{248, 248, 248}
\definecolor{coderule}{RGB}{210, 210, 210}
\definecolor{keywordblue}{RGB}{0, 100, 180}
\definecolor{commentgray}{RGB}{120, 120, 120}
\definecolor{stringgreen}{RGB}{30, 130, 30}

% Tie the Madrid structural colour to the book's link blue.
\setbeamercolor{structure}{fg=linkblue}
\setbeamercolor{title}{fg=white}
\setbeamercolor{frametitle}{fg=white}

% --- "the bigger picture" : narrative / intuition framing -----------
\newtcolorbox{bigpicture}{%
  enhanced, breakable, boxsep=2pt,
  colback  = narrativebg,
  colframe = narrativerule,
  boxrule  = 0pt, toprule = 1.4pt, bottomrule = 0.4pt,
  leftrule = 0pt, rightrule = 0pt, arc = 0pt,
  left = 6pt, right = 6pt, top = 4pt, bottom = 5pt,
  fonttitle = \footnotesize\scshape\color{narrativerule},
  title = {the bigger picture}, attach title to upper = {\par\smallskip},
}

% --- "under the hood" : the mechanism / technical detail ------------
\newtcolorbox{underhood}{%
  enhanced, breakable, boxsep=2pt,
  colback  = deepdivebg,
  colframe = linkblue,
  boxrule  = 0pt, toprule = 1.4pt, bottomrule = 0.4pt,
  leftrule = 0pt, rightrule = 0pt, arc = 0pt,
  left = 6pt, right = 6pt, top = 4pt, bottom = 5pt,
  fonttitle = \footnotesize\scshape\color{linkblue},
  title = {under the hood}, attach title to upper = {\par\smallskip},
}

% --- Python code style ---------------------------------------------
\lstset{
  basicstyle    = \ttfamily\footnotesize,
  keywordstyle  = \color{keywordblue}\bfseries,
  commentstyle  = \color{commentgray}\itshape,
  stringstyle   = \color{stringgreen},
  showstringspaces = false,
  language      = Python,
  breaklines    = true,
  backgroundcolor = \color{codebg},
  frame         = single,
  rulecolor     = \color{coderule},
  numbers       = none,
  aboveskip     = 4pt, belowskip = 4pt,
}

% --- Appendix conventions ------------------------------------------
% \appendixstart{Title} : begins the "extra material" appendix. Frames
% after it are not counted toward the main deck's frame total, keeping
% the core talk lean while preserving the deeper / optional content.
\newcommand{\appendixstart}[1]{%
  \appendix
  \setbeamertemplate{frametitle continuation}{}%
  \begin{frame}[noframenumbering]
    \begin{center}
      {\Large\scshape\color{linkblue} Appendix}\\[4pt]
      {\large #1}\\[10pt]
      {\footnotesize Extra / optional material — beyond the core lecture.}
    \end{center}
  \end{frame}
}
% Use \begin{frame}[noframenumbering]{...} for each appendix frame.

%   \title{...}\subtitle{...}\author{...}\institute{...}\date{\today}
%   \begin{document} ... \end{document}
% =====================================================================

\usetheme{Madrid}
\usecolortheme{whale}
\setbeamertemplate{navigation symbols}{}
\setbeamertemplate{caption}[numbered]
\setbeamercovered{transparent}

% --- Density controls: keep dense, equation-heavy frames on one slide ---
% Slightly smaller body text + tighter lists/blocks. Frame titles and the
% Madrid chrome keep their normal size, so the deck still reads as Madrid,
% just with more vertical room for content.
\setbeamerfont{normal text}{size=\small}
\setbeamertemplate{itemize/enumerate body begin}{\small}
\setbeamertemplate{itemize/enumerate subbody begin}{\footnotesize}
% Tighter block spacing.
\setbeamertemplate{block begin}{%
  \par\vskip2pt%
  \begin{beamercolorbox}[colsep*=.75ex,leftskip=1ex,rightskip=1ex]{block title}%
    \usebeamerfont*{block title}\insertblocktitle%
  \end{beamercolorbox}%
  {\parskip0pt\vskip-1pt}%
  \begin{beamercolorbox}[colsep*=.75ex,vmode,leftskip=1ex,rightskip=1ex]{block body}%
  \vskip1pt}
\setbeamertemplate{block end}{\end{beamercolorbox}\vskip2pt}

\usepackage{amsmath, amssymb}
\usepackage{booktabs}
\usepackage{xcolor}
\usepackage{listings}
\usepackage[skins, breakable]{tcolorbox}

% --- Book colour palette (kept in sync with book/preamble.tex) ---
\definecolor{linkblue}{RGB}{14, 75, 140}
\definecolor{deepdivebg}{RGB}{235, 244, 255}
\definecolor{narrativebg}{RGB}{255, 252, 240}
\definecolor{narrativerule}{RGB}{180, 130, 40}
\definecolor{steelblue}{RGB}{70, 130, 180}
\definecolor{forestgreen}{RGB}{34, 139, 34}
\definecolor{crimson}{RGB}{220, 20, 60}
\definecolor{darkorange}{RGB}{255, 140, 0}
\definecolor{codebg}{RGB}{248, 248, 248}
\definecolor{coderule}{RGB}{210, 210, 210}
\definecolor{keywordblue}{RGB}{0, 100, 180}
\definecolor{commentgray}{RGB}{120, 120, 120}
\definecolor{stringgreen}{RGB}{30, 130, 30}

% Tie the Madrid structural colour to the book's link blue.
\setbeamercolor{structure}{fg=linkblue}
\setbeamercolor{title}{fg=white}
\setbeamercolor{frametitle}{fg=white}

% --- "the bigger picture" : narrative / intuition framing -----------
\newtcolorbox{bigpicture}{%
  enhanced, breakable, boxsep=2pt,
  colback  = narrativebg,
  colframe = narrativerule,
  boxrule  = 0pt, toprule = 1.4pt, bottomrule = 0.4pt,
  leftrule = 0pt, rightrule = 0pt, arc = 0pt,
  left = 6pt, right = 6pt, top = 4pt, bottom = 5pt,
  fonttitle = \footnotesize\scshape\color{narrativerule},
  title = {the bigger picture}, attach title to upper = {\par\smallskip},
}

% --- "under the hood" : the mechanism / technical detail ------------
\newtcolorbox{underhood}{%
  enhanced, breakable, boxsep=2pt,
  colback  = deepdivebg,
  colframe = linkblue,
  boxrule  = 0pt, toprule = 1.4pt, bottomrule = 0.4pt,
  leftrule = 0pt, rightrule = 0pt, arc = 0pt,
  left = 6pt, right = 6pt, top = 4pt, bottom = 5pt,
  fonttitle = \footnotesize\scshape\color{linkblue},
  title = {under the hood}, attach title to upper = {\par\smallskip},
}

% --- Python code style ---------------------------------------------
\lstset{
  basicstyle    = \ttfamily\footnotesize,
  keywordstyle  = \color{keywordblue}\bfseries,
  commentstyle  = \color{commentgray}\itshape,
  stringstyle   = \color{stringgreen},
  showstringspaces = false,
  language      = Python,
  breaklines    = true,
  backgroundcolor = \color{codebg},
  frame         = single,
  rulecolor     = \color{coderule},
  numbers       = none,
  aboveskip     = 4pt, belowskip = 4pt,
}

% --- Appendix conventions ------------------------------------------
% \appendixstart{Title} : begins the "extra material" appendix. Frames
% after it are not counted toward the main deck's frame total, keeping
% the core talk lean while preserving the deeper / optional content.
\newcommand{\appendixstart}[1]{%
  \appendix
  \setbeamertemplate{frametitle continuation}{}%
  \begin{frame}[noframenumbering]
    \begin{center}
      {\Large\scshape\color{linkblue} Appendix}\\[4pt]
      {\large #1}\\[10pt]
      {\footnotesize Extra / optional material — beyond the core lecture.}
    \end{center}
  \end{frame}
}
% Use \begin{frame}[noframenumbering]{...} for each appendix frame.

%   \title{...}\subtitle{...}\author{...}\institute{...}\date{\today}
%   \begin{document} ... \end{document}
% =====================================================================

\usetheme{Madrid}
\usecolortheme{whale}
\setbeamertemplate{navigation symbols}{}
\setbeamertemplate{caption}[numbered]
\setbeamercovered{transparent}

% --- Density controls: keep dense, equation-heavy frames on one slide ---
% Slightly smaller body text + tighter lists/blocks. Frame titles and the
% Madrid chrome keep their normal size, so the deck still reads as Madrid,
% just with more vertical room for content.
\setbeamerfont{normal text}{size=\small}
\setbeamertemplate{itemize/enumerate body begin}{\small}
\setbeamertemplate{itemize/enumerate subbody begin}{\footnotesize}
% Tighter block spacing.
\setbeamertemplate{block begin}{%
  \par\vskip2pt%
  \begin{beamercolorbox}[colsep*=.75ex,leftskip=1ex,rightskip=1ex]{block title}%
    \usebeamerfont*{block title}\insertblocktitle%
  \end{beamercolorbox}%
  {\parskip0pt\vskip-1pt}%
  \begin{beamercolorbox}[colsep*=.75ex,vmode,leftskip=1ex,rightskip=1ex]{block body}%
  \vskip1pt}
\setbeamertemplate{block end}{\end{beamercolorbox}\vskip2pt}

\usepackage{amsmath, amssymb}
\usepackage{booktabs}
\usepackage{xcolor}
\usepackage{listings}
\usepackage[skins, breakable]{tcolorbox}

% --- Book colour palette (kept in sync with book/preamble.tex) ---
\definecolor{linkblue}{RGB}{14, 75, 140}
\definecolor{deepdivebg}{RGB}{235, 244, 255}
\definecolor{narrativebg}{RGB}{255, 252, 240}
\definecolor{narrativerule}{RGB}{180, 130, 40}
\definecolor{steelblue}{RGB}{70, 130, 180}
\definecolor{forestgreen}{RGB}{34, 139, 34}
\definecolor{crimson}{RGB}{220, 20, 60}
\definecolor{darkorange}{RGB}{255, 140, 0}
\definecolor{codebg}{RGB}{248, 248, 248}
\definecolor{coderule}{RGB}{210, 210, 210}
\definecolor{keywordblue}{RGB}{0, 100, 180}
\definecolor{commentgray}{RGB}{120, 120, 120}
\definecolor{stringgreen}{RGB}{30, 130, 30}

% Tie the Madrid structural colour to the book's link blue.
\setbeamercolor{structure}{fg=linkblue}
\setbeamercolor{title}{fg=white}
\setbeamercolor{frametitle}{fg=white}

% --- "the bigger picture" : narrative / intuition framing -----------
\newtcolorbox{bigpicture}{%
  enhanced, breakable, boxsep=2pt,
  colback  = narrativebg,
  colframe = narrativerule,
  boxrule  = 0pt, toprule = 1.4pt, bottomrule = 0.4pt,
  leftrule = 0pt, rightrule = 0pt, arc = 0pt,
  left = 6pt, right = 6pt, top = 4pt, bottom = 5pt,
  fonttitle = \footnotesize\scshape\color{narrativerule},
  title = {the bigger picture}, attach title to upper = {\par\smallskip},
}

% --- "under the hood" : the mechanism / technical detail ------------
\newtcolorbox{underhood}{%
  enhanced, breakable, boxsep=2pt,
  colback  = deepdivebg,
  colframe = linkblue,
  boxrule  = 0pt, toprule = 1.4pt, bottomrule = 0.4pt,
  leftrule = 0pt, rightrule = 0pt, arc = 0pt,
  left = 6pt, right = 6pt, top = 4pt, bottom = 5pt,
  fonttitle = \footnotesize\scshape\color{linkblue},
  title = {under the hood}, attach title to upper = {\par\smallskip},
}

% --- Python code style ---------------------------------------------
\lstset{
  basicstyle    = \ttfamily\footnotesize,
  keywordstyle  = \color{keywordblue}\bfseries,
  commentstyle  = \color{commentgray}\itshape,
  stringstyle   = \color{stringgreen},
  showstringspaces = false,
  language      = Python,
  breaklines    = true,
  backgroundcolor = \color{codebg},
  frame         = single,
  rulecolor     = \color{coderule},
  numbers       = none,
  aboveskip     = 4pt, belowskip = 4pt,
}

% --- Appendix conventions ------------------------------------------
% \appendixstart{Title} : begins the "extra material" appendix. Frames
% after it are not counted toward the main deck's frame total, keeping
% the core talk lean while preserving the deeper / optional content.
\newcommand{\appendixstart}[1]{%
  \appendix
  \setbeamertemplate{frametitle continuation}{}%
  \begin{frame}[noframenumbering]
    \begin{center}
      {\Large\scshape\color{linkblue} Appendix}\\[4pt]
      {\large #1}\\[10pt]
      {\footnotesize Extra / optional material — beyond the core lecture.}
    \end{center}
  \end{frame}
}
% Use \begin{frame}[noframenumbering]{...} for each appendix frame.


\title{Practical Session 09: Financial NLP and Sentiment}
\subtitle{Hands-On --- LM Scoring, FinBERT vs.\ Lexicon, and a Filing-Date Event Study}
\author{Juan F.\ Imbet}
\institute{Paris Dauphine -- PSL University}
\date{\today}

\begin{document}

\begin{frame}\titlepage\end{frame}

% ====================================================================
\begin{frame}{Session Overview}
\textbf{Duration:} 1 hour (50 min work + 10 min debrief)

\medskip
\textbf{Agenda:}
\begin{enumerate}
  \item \textbf{Recap} (5 min): LM net sentiment, macro-F1, CAR.
  \item \textbf{Problem 1} (12 min): Score an MD\&A excerpt with the LM lexicon by hand.
  \item \textbf{Problem 2} (12 min): Macro-F1 vs.\ accuracy on an imbalanced
        3-class confusion matrix.
  \item \textbf{Problem 3} (8 min): Build a $\mathrm{CAR}[1,60]$ predictive regression.
  \item \textbf{Case Study} (15 min): Run the chapter-09 companion notebook ---
        LM lexicon figure + net sentiment.
  \item \textbf{Discussion} (8 min): Lexicon vs.\ FinBERT vs.\ LLM in production.
\end{enumerate}

\medskip
\begin{block}{Code}
\texttt{code/notebooks/09-financial-nlp-sentiment/} ---
\texttt{demo.ipynb} (deterministic) and \texttt{exercises.ipynb} (live EDGAR lab).
\end{block}
\end{frame}

% ====================================================================
\begin{frame}{Recap: The Three Formulas You Need}
\textbf{LM net sentiment} (Problem 1):
\[
S_{\text{LM}}(d)=\frac{N_{\text{pos}}(d)-N_{\text{neg}}(d)}{N_{\text{words}}(d)}
\]

\textbf{Macro-F1} over 3 classes (Problem 2):
\[
F1_c=\frac{2\,\text{P}_c\text{R}_c}{\text{P}_c+\text{R}_c},\qquad
\text{macro-F1}=\tfrac13\sum_{c}F1_c
\]
with $\text{P}_c=\tfrac{TP_c}{TP_c+FP_c}$, $\text{R}_c=\tfrac{TP_c}{TP_c+FN_c}$.

\textbf{Cumulative abnormal return} (Problem 3):
\[
\mathrm{CAR}_i[1,60]=\sum_{\tau=1}^{60}\bigl(R_{i,\tau}-\hat R_{i,\tau}\bigr),
\quad
\mathrm{CAR}_i[1,60]=\alpha+\beta_1\mathrm{SUE}_{i,0}+\beta_2\tilde S_{i,0}+\varepsilon_i
\]
\end{frame}

% ====================================================================
\begin{frame}{Recap: Standardize Before You Compare}
Raw scores are not comparable across firms or time. Z-score within industry-year:
\[
\tilde S_{i,t}=\frac{S_{i,t}-\mu_{j(i),t}}{\sigma_{j(i),t}}
\]

\medskip
\begin{block}{Why it matters for the regression}
A defence or chemicals firm has structurally negative LM language. Without the
$\tilde S$ transform, $\beta_2$ picks up an industry fixed effect, not a sentiment
signal. The event study uses $\tilde S$, never the raw $S_{\text{LM}}$.
\end{block}

\medskip
\textbf{Also check:} ADF/KPSS for long series; one fixed model vintage to avoid
sentiment inflation; align timestamps to market hours (no look-ahead).
\end{frame}

% ====================================================================
\begin{frame}{Problem 1: LM Scoring by Hand}
\textbf{Excerpt} (10-K MD\&A):
\begin{exampleblock}{}
\textit{``Operating margins contracted as input costs rose. Management remains
uncertain about demand but is confident the restructuring will deliver
improved profitability.''}
\end{exampleblock}

\medskip
\textbf{Tasks.} Using the LM lists (negative dominates; \emph{uncertain} is on the
uncertainty list; \emph{cost} is \alert{neutral} in finance, not negative):
\begin{enumerate}
  \item Identify LM negative, positive, and uncertainty hits.
  \item Take $N_{\text{words}}=24$. Compute $S_{\text{LM}}$.
  \item A Harvard-GI classifier would also flag \emph{costs} and \emph{contracted}.
        Recompute the GI-style score. How large is the distortion?
\end{enumerate}

\medskip
\textbf{Discussion:} which word does the lexicon get \emph{wrong} on negation,
and how would FinBERT handle it?
\end{frame}

\begin{frame}{Solution: Problem 1}
\textbf{LM hits:}
\begin{itemize}\small
  \item Negative: \emph{contracted}, \emph{uncertain}(also uncertainty) $\to$ count
        \emph{contracted} as 1 negative.
  \item Positive: \emph{confident}, \emph{improved} $=2$.
  \item Uncertainty: \emph{uncertain} $=1$ (hedging signal, not polarity).
  \item \emph{costs}, \emph{rose} are \textbf{neutral} in LM.
\end{itemize}
\[
S_{\text{LM}}=\frac{2-1}{24}=+0.042 \quad(\text{mildly positive})
\]

\medskip
\textbf{GI-style:} adds \emph{costs} and \emph{contracted} as negative $\to$
negatives $=3$:
\[
S_{\text{GI}}=\frac{2-3}{24}=-0.042
\]
\alert{Sign flips} --- the GI false positives on neutral finance vocabulary reverse
the conclusion. This is exactly the Loughran--McDonald (2011) critique.

\medskip
\textbf{Negation:} ``not confident'' would still score \emph{positive} under any
lexicon; FinBERT's contextual \texttt{[CLS]} embedding resolves it.
\end{frame}

% ====================================================================
\begin{frame}{Problem 2: Macro-F1 vs.\ Accuracy}
A 3-class classifier on 1{,}000 financial sentences (neutral-heavy). Confusion
matrix (rows = true, cols = predicted):

\medskip
\begin{center}
\begin{tabular}{l|ccc|c}
\toprule
 & pred neg & pred neu & pred pos & total \\
\midrule
true neg & 60 & 35 & 5  & 100 \\
true neu & 20 & 560 & 20 & 600 \\
true pos & 5  & 90 & 205 & 300 \\
\bottomrule
\end{tabular}
\end{center}

\medskip
\textbf{Tasks.}
\begin{enumerate}
  \item Compute overall accuracy.
  \item Compute per-class precision, recall, F1 for \textbf{negative}.
  \item Given $F1_{\text{neu}}\approx 0.86$ and $F1_{\text{pos}}\approx 0.74$,
        compute macro-F1.
  \item Why is the gap between accuracy and macro-F1 the point of this exercise?
\end{enumerate}
\end{frame}

\begin{frame}{Solution: Problem 2}
\textbf{Accuracy} $=(60+560+205)/1000 = \mathbf{0.825}$.

\medskip
\textbf{Negative class:}
\begin{itemize}\small
  \item $TP=60$, $FP=20+5=25$, $FN=35+5=40$.
  \item Precision $=60/85 = 0.706$; \quad Recall $=60/100 = 0.600$.
  \item $F1_{\text{neg}} = \dfrac{2(0.706)(0.600)}{0.706+0.600} = \mathbf{0.649}$.
\end{itemize}

\medskip
\textbf{Macro-F1} $=\tfrac13(0.649+0.86+0.74) = \mathbf{0.750}$.

\medskip
\begin{alertblock}{The lesson}
Accuracy 0.825 looks strong, but it is propped up by the 600 easy neutral cases.
Macro-F1 0.750 reveals the model is markedly weaker on the economically important
\textbf{negative} class ($F1=0.65$). For a trading or risk signal, the
neg/pos tails are what you act on --- so report macro-F1.
\end{alertblock}
\end{frame}

% ====================================================================
\begin{frame}{Problem 3: Designing the Event Study}
\textbf{Setup.} 800 earnings calls. For each you have: the call-day standardized
text sentiment $\tilde S_{i,0}$, the earnings surprise
$\mathrm{SUE}_{i,0}=(\mathrm{EPS}-\widehat{\mathrm{EPS}})/\hat\sigma_i$, and daily returns.

\medskip
\textbf{Tasks.}
\begin{enumerate}
  \item Write $\mathrm{CAR}_i[1,60]$ in terms of abnormal returns. What benchmark
        $\hat R$ would you use, and why not raw returns?
  \item State the regression and the null/alternative on $\beta_2$.
  \item You find $\hat\beta_2 = 0.018$ ($t=2.6$) after controlling for SUE.
        Interpret economically.
  \item Name two robustness checks before you trust this.
\end{enumerate}
\end{frame}

\begin{frame}{Solution: Problem 3}
\textbf{1.} $\mathrm{CAR}_i[1,60]=\sum_{\tau=1}^{60}(R_{i,\tau}-\hat R_{i,\tau})$,
$\hat R$ from a market model / CAPM / factor benchmark. Raw returns mix in
market-wide moves; abnormal returns isolate the firm-specific drift.

\medskip
\textbf{2.} $\mathrm{CAR}_i[1,60]=\alpha+\beta_1\mathrm{SUE}_{i,0}+\beta_2\tilde S_{i,0}+\varepsilon_i$.
$H_0:\beta_2=0$ (text adds nothing beyond the EPS number);
$H_1:\beta_2>0$ (residual text sentiment predicts drift).

\medskip
\textbf{3.} $\hat\beta_2=0.018$, $t=2.6$: a one-SD increase in standardized
sentiment $\to$ +1.8 pp of 60-day abnormal return, \emph{orthogonal} to the
quantitative surprise. Text carries incremental information --- post-announcement
drift conditioned on tone.

\medskip
\textbf{4.} Robustness: (i) strictly out-of-sample period not used for model
selection (economic-validity test); (ii) cluster SEs by firm/date; also check
look-ahead (timestamp vs.\ market close) and a placebo pre-event window.
\end{frame}

% ====================================================================
\begin{frame}[fragile]{Case Study: Chapter-09 Companion Notebook}
Open \texttt{code/notebooks/09-financial-nlp-sentiment/demo.ipynb}. It is
deterministic (no network) and reproduces the LM lexicon figure + the worked
net-sentiment example.

\begin{lstlisting}
import matplotlib.pyplot as plt

# Loughran-McDonald finance dictionary word-list sizes.
lists = {"Negative": 2355, "Litigious": 903, "Positive": 354,
         "Uncertainty": 297, "Weak modal": 27, "Strong modal": 19}
items = sorted(lists.items(), key=lambda kv: kv[1])
plt.barh([k for k, _ in items], [v for _, v in items], color="#4c72b0")
plt.xlabel("Number of words"); plt.title("Loughran-McDonald dictionary")
plt.tight_layout(); plt.show()
\end{lstlisting}
\end{frame}

\begin{frame}[fragile]{Case Study: Net LM Sentiment Example}
The same notebook reproduces the worked net-sentiment number from the chapter.

\begin{lstlisting}
# Net LM sentiment for the chapter MD&A excerpt (ex:lm-example).
n_pos, n_neg, n_words = 1, 4, 35
s_lm = (n_pos - n_neg) / n_words
print(f"S_LM = ({n_pos} - {n_neg}) / {n_words} = {s_lm:.3f}")  # ~ -0.086
\end{lstlisting}
\end{frame}

\begin{frame}{Case Study: Diagnostic Questions}
Run \texttt{demo.ipynb}, then answer:

\medskip
\textbf{Q1.} The negative list (2{,}355) dwarfs the positive list (354). What does
this asymmetry tell you about the language of corporate disclosure --- and about
the \emph{base rate} a lexicon will produce on a typical 10-K?

\medskip
\textbf{Q2.} The excerpt scores $S_{\text{LM}}=-0.086$. Re-run with
\texttt{n\_pos, n\_neg = 2, 2}. How sensitive is the sign to a single word-list
membership decision? What does this imply for reproducibility?

\medskip
\textbf{Q3 (live lab).} Move to \texttt{exercises.ipynb}: it pulls real 10-Ks from
EDGAR and prices from \texttt{yfinance}. Score one filing, then test whether
filing-day LM negativity is associated with the filing-day abnormal return.
Does the sign match Loughran--McDonald (2011)?

\medskip
\begin{block}{Extension}
Swap the lexicon for \texttt{yiyanghkust/finbert-tone} on the same filings.
Where do the two methods \emph{disagree}, and which sentences drive it?
\end{block}
\end{frame}

% ====================================================================
\begin{frame}{Discussion: Choosing a Method in Production}
For each scenario, pick lexicon / FinBERT / LLM and justify on cost, latency,
labels, interpretability, and accuracy.

\medskip
\begin{enumerate}
  \item A regulator must audit \emph{why} every sentiment score on a 10-K was
        assigned. 50k filings/year, no labelled data, full transparency required.
  \item A desk scores 5M news headlines/day with a $<50$ ms/headline budget and
        10k labelled examples in hand.
  \item A research team prototypes sentiment on a brand-new document type
        (crypto whitepapers) with \emph{zero} labels and a small corpus.
\end{enumerate}

\medskip
\begin{block}{Frame your answer around}
Auditability (lexicon wins), throughput/latency (distilled FinBERT wins),
cold-start flexibility (LLM zero-shot wins) --- and when you would \emph{distil}
an LLM into a small model for scale.
\end{block}
\end{frame}

\begin{frame}{Wrap-Up: Key Takeaways}
\begin{itemize}
  \item \textbf{Lexicons are auditable but brittle:} the GI vs.\ LM sign flip in
        Problem 1 is the whole case for domain-specific dictionaries.
  \item \textbf{Report macro-F1, not accuracy:} the neutral majority hides weakness
        on the actionable tails (Problem 2).
  \item \textbf{Standardize, then test economically:} z-score within industry-year,
        run the CAR regression, validate strictly out-of-sample (Problem 3).
  \item \textbf{Method choice is a trade-off:} auditability vs.\ throughput vs.\
        cold-start flexibility --- no universal winner.
\end{itemize}

\medskip
\begin{block}{Next practical}
Practical 10 --- domain-specific LLMs: fine-tuning and adaptation of FinBERT's
successors on long financial documents.
\end{block}
\end{frame}

% ====================================================================
% APPENDIX
% ====================================================================
\appendixstart{Stretch Problems}

\begin{frame}[noframenumbering]{Stretch: Krippendorff's Alpha}
Two annotators label 5 sentences (P=pos, N=neu, G=neg):
\begin{center}
\begin{tabular}{lccccc}
\toprule
        & s1 & s2 & s3 & s4 & s5 \\
\midrule
Annot.\ A & P & N & G & N & P \\
Annot.\ B & P & N & G & G & P \\
\bottomrule
\end{tabular}
\end{center}
\textbf{Tasks.} (1) Compute observed disagreement $D_o$ (fraction of disagreeing
pairs). (2) With marginal frequencies, compute $D_e=1-\sum_c(n_c/n)^2$.
(3) Compute $\alpha=1-D_o/D_e$. Is it above Krippendorff's 0.667 reliability bar?
(4) Why would the \emph{ordinal} variant give a different, more lenient verdict
on the s4 P/G... here N/G disagreement?
\end{frame}

\begin{frame}[noframenumbering]{Stretch: From Score to Sharpe}
\textbf{Setup.} A standardized news-sentiment signal $\tilde S$ is rebalanced
daily into a long--short portfolio (long top decile, short bottom decile).

\medskip
\textbf{Tasks.}
\begin{enumerate}\small
  \item The in-sample annualized Sharpe is 1.4; strictly out-of-sample it is 0.5.
        Name three reasons for the decay (hint: one is in the lecture's
        ``sentiment inflation'' box).
  \item Kirtac \& Germano (2024) report OPT beats the LM dictionary on long--short
        Sharpe. Why might the gap shrink \emph{after} transaction costs at high turnover?
  \item Lehner (2024) shows LLMs can rewrite filings to shift measured tone.
        How does adversarial text-generation threaten a filings-based signal's
        out-of-sample stability?
\end{enumerate}
\end{frame}

\end{document}
